\chapter{Evaluation}
\label{ch:-evaluation}

This chapter evaluates the performance of the drug-target interaction (DTI) prediction model across multiple iterations of development.
The evaluation will be structured around three key areas: \textbf{1}: model progression through iterative improvements (training, architecture, framework); \textbf{2}: sensitivity to data and decision thresholds; and \textbf{3}: comparison of the final model against established state-of-the-art DTI methods on BindingDB.
Through this comprehensive evaluation, we aim to demonstrate the effectiveness of our modelling choices and identify areas for future improvement.

% ============================================================
\section{Evaluation Methodology}
\label{sec: eval_methodology}
% (~200 words)
% Justify your choices BEFORE presenting results.
This section will outline the components of the evaluation methodology, including the metrics used to assess model performance, the data splits for training and testing, and the rationale behind the choice of a decision threshold for classification.
    \subsection{Metrics}
    \label{sec: eval_metrics}
    The following metrics are computed on the training, validation, and test set predictions:    \begin{itemize}
        \item Area Under the Receiver Operating Characteristic Curve (AUC-ROC)
        \item Area Under the Precision-Recall Curve (PR-AUC)
        \item F1 Score
        \item Precision
        \item Recall (Sensitivity)
        \item Accuracy
        \item Confusion Matrix
        \item Specificity
    \end{itemize}

Accuracy is a fundamental metric that measures the performance by finding out the ratio between correct predictions and all predictions.
This alone would be meaningless because in an imbalanced dataset, a model could achieve high accuracy by simply predicting the majority class.
Precision and recall are crucial in the context of DTI prediction, where false positives (predicting a drug-target interaction that does not exist) can lead to wasted resources in experimental validation, while false negatives (failing to predict a true interaction) can result in missed opportunities for drug discovery.
The F1 score balances precision and recall into a single metric, providing a more holistic view of model performance.
AUC-ROC and PR-AUC are threshold-independent metrics that summarise model performance across all possible decision thresholds.
AUC-ROC captures the trade-off between the true positive rate and the false positive rate, and can be interpreted as the probability that the model assigns a higher score to a random positive instance than to a random negative one.
PR-AUC instead captures the trade-off between precision and recall, which is particularly informative with class imbalance.
Finally, the confusion matrix provides a detailed breakdown of true positives, false positives, true negatives, and false negatives, allowing for a granular analysis of model performance \citep{geeksforgeeks_metrics_2025}.

    \subsection{Data Splits and Reproducibility}
    \label{sec: eval_splits}
This work uses a fixed 65/20/15 split of the BindingDB data for the training, validation, and test sets.
However, this is contrasting a 10-fold cross-validation approach used by DeepCDA and MINDG, which means that the reported metrics from those papers are not directly comparable to the results, but rather serve as approximate reference points.
The data splits are also varied to assess their effect on performance.
    \subsection{Decision Threshold}
    \label{sec: eval_threshold}
Another important aspect of the evaluation is the choice of the decision threshold for classifying predictions as positive or negative interactions.
We have opted for pAffinity value $> 7.0$ determining binder and pAffinity value $< 5.3$ determining non-binder with a grey area of 1.70 pAffinity units as mentioned in~\ref{sec: data_threshold}.
However, to improve the robustness of the model, the setup is experimented with to fully evaluate the model.

% ============================================================
\section{Model Evaluation}
\label{sec:model_eval}
Evaluation of the model opens with a basic NumPy implementation and iteratively improving it through changes to the loss function, architecture, and training procedure.

    % --------------------------------------------------------

    \subsection{Experimental Setup}
    \label{sec: eval_setup}
All models were run on the same hardware (MacBook Air M3, 16\,GB RAM) to ensure a fair comparison of training times and computational efficiency.
The NumPy iterations were trained with a maximum of 1000 epochs and the PyTorch model for 100, both will introduce and implement early stopping based on validation loss to prevent overfitting and reduce unnecessary computation.
We cap the NumPy to 1000 epochs to allow for sufficient convergence under full-batch gradient descent; the PyTorch model was restricted to a tenth of that because of the mini-batch gradient descent and to reduce the chance of overfitting.
This provided the PyTorch model with more gradient updates per epoch, so equivalent optimisation progress is made in far fewer epochs.
To keep consistency throughout, each iteration will be evaluated on the same subset of the BindingDB dataset, the class distribution of which is shown in Table~\ref{tab:class-distribution}.
This 18,520-row subset has already had the grey area ($5.3 \leq$ \textit{pAffinity} $\leq 7.0$) removed during the preprocessing stage described in Chapter~\ref{chap:data-pipeline}, so the iterative experiments are trained and evaluated only on high-confidence binders and non-binders.
\begin{table}[H]
\centering
\caption{Class distribution across train, validation, and test splits.}
\label{tab:class-distribution}
\begin{tabular}{lrrr}
\toprule
\textbf{Split} & \textbf{Binders} & \textbf{Non-binders} & \textbf{Total} \\
\midrule
Train      & 7,291  & 4,747 & 12,038 \\
Validation & 2,226  & 1,478 & 3,704  \\
Test       & 1,671  & 1,107 & 2,778  \\
\midrule
Total      & 11,188 & 7,332 & 18,520 \\
\bottomrule
\end{tabular}
\end{table}

    \subsection{Baseline: Basic NumPy Model}
    \label{sec: eval_basic}
    % (~200 words + figures)
    To establish a baseline for the DTI prediction task, we implemented a simple feedforward neural network using only NumPy and a deliberately minimal baseline with a low complexity.
    This model consists of a single hidden layer with sigmoid activations, which is also present on the output layer, and it is trained using Binary Cross-Entropy (BCE) loss.
    The absence of many components intentionally allows for incremental improvement, and any improvement in performance can be attributed to a specific change and not a combination of them.
    Table~\ref{tab:hyperparams} showcases the setup for the initial model.
  \begin{table}[H]
  \centering
  \caption{Hyperparameter selection in the NumPy baseline.
           Dashes (\textemdash) indicate features absent from the baseline model.}
  \label{tab:hyperparams}
  \begin{tabular}{lc}
  \toprule
  \textbf{Parameter} & \textbf{Basic (NumPy)} \\
  \midrule
  \textit{Architecture} \\
  \quad Hidden size          & 32 \\
  \quad Activation           & Sigmoid \\
  \quad Batch normalisation  & \textemdash \\
  \quad Dropout              & \textemdash \\
  \midrule
  \textit{Training} \\
  \quad Loss function        & BCE \\
  \quad Optimiser            & \textemdash \\
  \quad Learning rate        & 0.5 \\
  \quad Weight decay         & \textemdash \\
  \quad Batch size           & Full batch \\
  \quad Max epochs           & 1000 \\
  \midrule
  \textit{Regularisation} \\
  \quad Early stopping       & \textemdash \\
  \quad LR scheduler         & \textemdash \\
  \bottomrule
  \end{tabular}
  \end{table}

After training this baseline model, we evaluated its performance against the validation and test sets.
The results of key metrics are summarised in the following graph:
\begin{figure}[H]
    \centering
    \includegraphics[width=0.8\textwidth]{/Users/drs/Projects/DTI/Backend/app/Models/basic/figures/basic_BCE_metrics_table.png}
    \caption{Baseline model performance across key metrics.}
    \label{fig:basic_metrics_table}
\end{figure}
The data in this table shows strong consistency across splits, which is a sign of minimal overfitting and deals well with unseen data.
A recall of $\sim$0.76 indicates the model is good at catching positive cases, which is important when predicting a false negative is costly.
A strong recall is important for drug discovery, as missing a true drug-target interaction could mean overlooking a potential binding candidate.
The AUC-ROC of $\sim$0.75 suggests the model has a reasonable ability to distinguish between binders and non-binders, although there is certainly room for improvement.
The Precision score of $\sim$0.75 indicates that when the model predicts a positive interaction, it is correct ~75\% of the time, which is decent but could lead to some false positives that would require experimental validation.

    \subsubsection{Hyperparameter tuning}
To further evaluate the baseline model, a hyperparameter sweep was conducted by varying the learning rate and hidden layer size.
The Figure~\ref{fig:hyperparam_sweep} visualises the sweep of learning rate and hidden size visualising the difference in each parameter:
\begin{figure}[H]
    \centering
    \begin{minipage}{0.48\textwidth}
        \centering
        \includegraphics[width=\textwidth]{/Users/drs/Projects/DTI/Backend/app/Models/basic/figures/sweep_lr_loss_curve.png}
        \caption*{Learning Rate Sweep}
    \end{minipage}
    \hfill
    \begin{minipage}{0.48\textwidth}
        \centering
        \includegraphics[width=\textwidth]{/Users/drs/Projects/DTI/Backend/app/Models/basic/figures/sweep_hidden_metrics.png}
        \caption*{Hidden Layer Size Sweep}
    \end{minipage}
    \caption{Hyperparameter tuning results for the baseline model.}
    \label{fig:hyperparam_sweep}
\end{figure}
From the learning rate sweep we learned that 0.1 achieves the same result as 0.5 but with a smoother loss curve, which is a sign of more stable training.
From the hidden layer size sweep we learned that accuracy and F1 score are stable across all hidden sizes and AUC-ROC only improves marginally.
This suggests the basic model has reached a boundary in performance imposed by its single layer-architecture, and that further improvements will likely require architectural changes rather than just tuning hyperparameters.

    \subsection{Training Procedure advancements}
    \label{sec: eval_iter1}
Building on the basic model, progressive changes to the training procedure are tested to improve performance.
        \subsubsection{Minibatch Gradient Descent}
Minibatch gradient descent provides more stable and efficient updates compared to full batch training, especially on larger datasets.
The results of which are seen in Figure~\ref{fig:iter1_metrics_table}.
\begin{figure}[H]
    \centering
    \includegraphics[width=0.8\textwidth]{/Users/drs/Projects/DTI/Backend/app/Models/Training/figures/iter1_metrics_table.png}
    \caption{Iteration 1 model performance across key metrics.}
    \label{fig:iter1_metrics_table}
\end{figure}
The addition of mini-batch gradient descent improved recall substantially at the cost of a small drop in precision, reflecting the noisier gradient updates – a known trade-off that the subsequent iterations will aim to address.
And a different loss curve compared to the baseline, this change is because the model is now updating its weights more frequently, which can lead to faster convergence and a smoother loss curve compared to the baseline's full batch training.
\begin{figure}[H]
    \centering
    \begin{minipage}{0.48\textwidth}
        \centering
        \includegraphics[width=\textwidth]{/Users/drs/Projects/DTI/Backend/app/Models/basic/figures/basic_loss_curve.png}
        \caption*{Baseline (Full-batch GD)}
    \end{minipage}
    \hfill
    \begin{minipage}{0.48\textwidth}
        \centering
        \includegraphics[width=\textwidth]{/Users/drs/Projects/DTI/Backend/app/Models/Training/figures/iter1_loss_curve.png}
        \caption*{Iteration 1 (Mini-batch GD)}
    \end{minipage}
    \caption{Training loss curves comparing the baseline to Iteration 1.}
    \label{fig:loss_curve_comparison}
\end{figure}
The training of the baseline model, produced a significant instability spike.
This is a characteristic of full-batch gradient descent with an aggressive learning rate.
Mini-batch gradient descent eliminates this instability by providing more frequent updates, which allows the model to adjust its weights more smoothly and avoid large jumps in the loss landscape.
This results in a more stable training process, as evidenced by the smoother loss curve in Iteration 1 compared to the baseline.

    \subsubsection{He initialisation}
He initialisation was added alongside minibatch gradient descent to break the symmetry and prevent vanishing gradients at the start of training~\citep{he2015delving}.
He initialisation works by setting the initial weights of the network using this equation:
\begin{equation}
    W \sim \mathcal{N}\left(0, \sqrt{\frac{2}{n_{\text{in}}}}\right)
\end{equation}
where $n_{\text{in}}$ is the number of input units in the layer.
This method is particularly effective for layers with ReLU activations, as it helps maintain the variance of the activations throughout the network, which can lead to faster convergence and improved performance.
He initialisation was designed for ReLU activations specifically; applying it alongside sigmoid hidden units is a deliberate intermediate step, as ReLU hidden activations are introduced in the following iteration.
Nonetheless, the variance-preserving property of He initialisation still offers an improvement over naive small weight initialisation ($\times 0.01$) regardless of activation function.
These small improvements can be seen within Figure~\ref{fig:iter2_improvements}.
\begin{figure}[H]
    \centering
    \begin{minipage}{0.48\textwidth}
        \centering
        \includegraphics[width=\textwidth]{/Users/drs/Projects/DTI/Backend/app/Models/Training/Iterations/figures/iter2_metrics_table.png}
        \caption*{Iteration 2 model performance across key metrics.}
    \end{minipage}
    \hfill
    \begin{minipage}{0.48\textwidth}
        \centering
        \includegraphics[width=\textwidth]{/Users/drs/Projects/DTI/Backend/app/Models/Training/Iterations/figures/iter2_loss_curve.png}
        \caption*{Training loss curves comparing Iteration 1 to Iteration 2.}
    \end{minipage}
    \caption{Iteration 2 improvements.}
    \label{fig:iter2_improvements}
\end{figure}
The improvements from Iteration 1 to Iteration 2 are modest but consistent across all metrics.

    \subsubsection{ReLU activations}
To further improve the training stability and representational capacity of the model, sigmoid activations are replaced in the hidden layer with ReLU (Rectified Linear Unit) activations.
ReLU is defined as:
\begin{equation}
    \text{ReLU}(x) = \max(0, x)
\end{equation}
ReLU addresses the vanishing gradient problem associated with sigmoid activations, as it allows for a constant gradient for positive inputs, which helps maintain the flow of gradients during backpropagation, especially in deeper networks \citep{nair2010rectified}.
This change was impactful and can be seen in the Figure~\ref{fig:iter3_improvements}.
\begin{figure}[H]
    \centering
    \begin{minipage}{0.48\textwidth}
        \centering
        \includegraphics[width=\textwidth]{/Users/drs/Projects/DTI/Backend/app/Models/Training/Iterations/figures/iter3_metrics_table.png}
        \caption*{Iteration 3 model performance across key metrics.}
    \end{minipage}
    \hfill
    \begin{minipage}{0.48\textwidth}
        \centering
        \includegraphics[width=\textwidth]{/Users/drs/Projects/DTI/Backend/app/Models/Training/Iterations/figures/iter3_roc_curve.png}
        \caption*{AUC-ROC curves comparing Iteration 2 to Iteration 3.}
    \end{minipage}
    \caption{Iteration 3 improvements.}
    \label{fig:iter3_improvements}
\end{figure}
The switch to ReLU activations in Iteration 3 led to a noticeable improvement in all performance metrics, particularly in AUC-ROC and F1 score, which suggests that the model is now better able to capture complex patterns in the data and distinguish between binders and non-binders more effectively.

\subsubsection{Early Stopping}
The final iteration to the training procedure is to add early stopping based on validation loss, which can help prevent overfitting and reduce training time by halting training once the model's performance on the validation set starts to degrade.
After adding early stopping, the performance changes can be seen in Figure~\ref{fig:iter4_improvements}.
\begin{figure}[H]
    \centering
    \begin{minipage}{0.48\textwidth}
        \centering
        \includegraphics[width=\textwidth]{/Users/drs/Projects/DTI/Backend/app/Models/Training/Iterations/figures/iter4_metrics_table.png}
        \caption*{Iteration 4 model performance across key metrics.}
    \end{minipage}
    \hfill
    \begin{minipage}{0.48\textwidth}
        \centering
        \includegraphics[width=\textwidth]{/Users/drs/Projects/DTI/Backend/app/Models/Training/Iterations/figures/iter4_loss_curve.png}
        \caption*{Training loss curves showing early stopping in Iteration 4.}
    \end{minipage}
    \caption{Iteration 4 improvements with early stopping.}
    \label{fig:iter4_improvements}
\end{figure}
As seen in Figure~\ref{fig:iter4_improvements}, early stopping led to premature halting of training, missing out on some of the later epochs where the model could have achieved better performance.
This is a common issue with early stopping, as it relies on the validation loss to determine when to stop training, and if the validation loss is noisy or has a local minimum, it can cause the training to stop too early before the model has fully converged.
For these reasons, early stopping will not be used in the final model for validation loss but only for validation AUC.
However, for the rest of the iterative improvements we will be using iteration 3 as the basis for future advancements to architecture.

    \subsection{Architecture Advancements}
    \label{sec: eval_iter2}
Next, architectural changes to the model occur to further improve its performance.
\subsubsection{More Layers}
A second hidden layer was added to increase the representational capacity of the network.
This allows the model to learn more complex patterns in the data, which should lead to improved performance.
The effect of a second hidden layer can be seen in Figure~\ref{fig:arch1_improvements}.
\begin{figure}[H]
    \centering
    \begin{minipage}{0.48\textwidth}
        \centering
        \includegraphics[width=\textwidth]{/Users/drs/Projects/DTI/Backend/app/Models/Architecture/figures/arch1_metrics_table.png}
        \caption*{Architectural Iteration 1}
    \end{minipage}
    \hfill
    \begin{minipage}{0.48\textwidth}
        \centering
        \includegraphics[width=\textwidth]{/Users/drs/Projects/DTI/Backend/app/Models/Architecture/figures/arch1_loss_curve.png}
        \caption*{Loss curves comparing architecture Iteration 1 to training iteration 3.}
    \end{minipage}
    \caption{Improvements from adding a second hidden layer.}
    \label{fig:arch1_improvements}
\end{figure}
The addition of a second hidden layer unexpectedly led to a drop in performance across all metrics, which suggests that the model may be overfitting to the training data or that the additional layer is not being effectively trained with the current settings.
This highlights the importance of careful architectural design and the need for regularisation techniques to prevent overfitting when increasing model complexity.
Regularisation techniques such as dropout and batch normalisation will be explored in the next iteration to address these issues.

\subsubsection{Dropout Added}
To combat the overfitting created by a second hidden layer, dropout regularisation is added, which randomly zeroes out a fraction of the activations during training to prevent co-adaptation of neurons and encourage the model to learn more robust features \citep{srivastava2014dropout}.
We chose a low dropout rate of 0.2 compared with~\citet{srivastava2014dropout}, who recommend a dropout of 0.5 in large networks, to reflect the smaller scale of this model and the presence of other regularisation techniques.
The effect of dropout is shown in Figure~\ref{fig:arch2_improvements.
\begin{figure}[H]
    \centering
    \begin{minipage}{0.48\textwidth}
        \centering
        \includegraphics[width=\textwidth]{/Users/drs/Projects/DTI/Backend/app/Models/Architecture/figures/arch2_metrics_table.png}
        \caption*{Architectural Iteration 2}
    \end{minipage}
    \hfill
    \begin{minipage}{0.48\textwidth}
        \centering
        \includegraphics[width=\textwidth]{/Users/drs/Projects/DTI/Backend/app/Models/Architecture/figures/arch2_loss_curve.png}
        \caption*{Loss curves comparing architecture Iteration 2 to architecture Iteration 1.}
    \end{minipage}
    \caption{Improvements from adding dropout regularisation.}
    \label{fig:arch2_improvements}
\end{figure}
The addition of dropout led to a small improvement in precision while sacrificing some recall, which is a common trade-off when adding regularisation.
This suggests that dropout is helping to reduce overfitting by encouraging the model to learn more general features, but it may also be causing the model to miss some true positive interactions, which is reflected in the drop in recall.

\subsubsection{L2 Regularisation}
Because dropout alone hasn't fully addressed the overfitting issue, L2 regularisation (weight decay) was also added, which penalises large weights directly by adding a term to the loss function~\citep{goodfellow2016deep}:
\begin{equation}
    \mathcal{L}_{\text{total}} = \mathcal{L}_{\text{BCE}} +
    \frac{\lambda}{2} \sum_{w} w^2
\end{equation}
Where $\lambda$ is the regularisation strength and $w$ are the model weights.
After adding L2 regularisation of 0.01 alongside a dropout of 0.2, the change performance is shown in Figure~\ref{fig:arch3_improvements}.
\begin{figure}[H]
    \centering
    \begin{minipage}{0.48\textwidth}
        \centering
        \includegraphics[width=\textwidth]{/Users/drs/Projects/DTI/Backend/app/Models/Architecture/figures/arch3_metrics_table.png}
        \caption*{Architectural Iteration 3}
    \end{minipage}
    \hfill
    \begin{minipage}{0.48\textwidth}
        \centering
        \includegraphics[width=\textwidth]{/Users/drs/Projects/DTI/Backend/app/Models/Architecture/figures/arch3_loss_curve.png}
        \caption*{Loss curves comparing architecture Iteration 3 to architecture Iteration 2.}
    \end{minipage}
    \caption{Improvements from adding L2 regularisation.}
    \label{fig:arch3_improvements}
\end{figure}
Figure~\ref{fig:arch3_improvements} shows the addition of L2 regularisation led to a further improvement in precision, while recall dropped significantly.
This is visualised in the loss curve where the training loss is higher but the learning is more stable, which is a sign that the model is now learning more general features and is less likely to overfit to the training data.

To improve this model further, dropout is increased to 0.4 to counter the low recall and we want to make sure we are catching as many true positives as possible.
The result can be observed in Figure~\ref{fig:arch4_improvements}.
\begin{figure}[H]
    \centering
    \begin{minipage}{0.48\textwidth}
        \centering
        \includegraphics[width=\textwidth]{/Users/drs/Projects/DTI/Backend/app/Models/Architecture/figures/arch3_metrics_table(0.4).png}
        \caption*{Architectural Iteration 4}
    \end{minipage}
    \hfill
    \begin{minipage}{0.48\textwidth}
        \centering
        \includegraphics[width=\textwidth]{/Users/drs/Projects/DTI/Backend/app/Models/Architecture/figures/arch3_loss_curve(0.4).png}
        \caption*{Loss curves comparing architecture Iteration 4 to architecture Iteration 3.}
    \end{minipage}
    \caption{Improvements from increasing dropout to 0.4.}
    \label{fig:arch4_improvements}
\end{figure}
Figure~\ref{fig:arch4_improvements} shows that increasing the dropout to 0.4 led to a rebound of recall, which is a sign that the model is now able to capture more true positive interactions.
However, this came at the cost of a significant drop in precision, which suggests that the model is now making more false positive predictions.
This highlights the delicate balance between precision and recall when tuning regularisation parameters.


    \subsection{New framework}
    \label{sec: eval_iter3}
To fully address the overfitting issue, more significant changes to the training procedure and architecture take place.
\subsubsection{PyTorch Implementation}
The model in PyTorch is implemented with the same architecture and training procedure used in architectural iteration 4.
This is to maintain comparability whilst also leveraging the benefits of modern deep learning frameworks such as automatic differentiation, efficient GPU computation, and built-in regularisation techniques \citep{paszke2019pytorch}.
The initial implementation builds upon architectural iteration 4~\ref{fig:arch4_improvements} whilst also utilising a dual-branch PyTorch encoder to process drug and target features separately before concatenating them for the final prediction.
This had significant improvements shown in Figure~\ref{fig:pytorch_improvements}.
\begin{figure}[H]
    \centering
    \begin{minipage}{0.48\textwidth}
        \centering
        \includegraphics[width=\textwidth]{/Users/drs/Projects/DTI/Backend/app/Models/Torch/figures/dnn_sgd_metrics_table.png}
        \caption*{PyTorch DNN model performance across key metrics.}
    \end{minipage}
    \hfill
    \begin{minipage}{0.48\textwidth}
        \centering
        \includegraphics[width=\textwidth]{/Users/drs/Projects/DTI/Backend/app/Models/Torch/figures/dnn_sgd_roc_curve.png}
        \caption*{AUC-ROC curve for the PyTorch DNN model.}
    \end{minipage}
    \caption{Performance of the PyTorch DNN model.}
    \label{fig:pytorch_improvements}
\end{figure}

Seeing such improvement in the ROC curve and metrics table indicates that the PyTorch implementation is able to learn more complex patterns in the data and generalise better to unseen examples, which is likely due to the combination of architectural improvements and the benefits of using a modern deep learning framework.
However, this is at the expense of Recall, which suggests that the model is now making more false negative predictions, which is a concern in the context of DTI prediction where missing a true interaction could have significant implications for drug discovery.
This issue is addressed within the next iteration by including the Adam optimiser.

    \subsubsection{Adam Optimiser}
The Adam optimiser was adopted as the final optimisation strategy.
The improvements from using Adam can be seen in the following Figure~\ref{fig:adam_improvements}.
\begin{figure}[H]
    \centering
    \begin{minipage}{0.48\textwidth}
        \centering
        \includegraphics[width=\textwidth]{/Users/drs/Projects/DTI/Backend/app/Models/Torch/figures/dnn_adam_metrics_table.png}
        \caption*{PyTorch DNN with Adam optimiser performance across key metrics.}
    \end{minipage}
    \hfill
    \begin{minipage}{0.48\textwidth}
        \centering
        \includegraphics[width=\textwidth]{/Users/drs/Projects/DTI/Backend/app/Models/Torch/figures/dnn_adam_loss_curve.png}
        \caption*{Training loss curve for the PyTorch DNN with Adam optimiser.}
    \end{minipage}
    \caption{Performance improvements from using the Adam optimiser.}
    \label{fig:adam_improvements}
\end{figure}
The use of the Adam optimiser led to a significant improvement in recall without sacrificing too much, which is a sign that the model is now able to capture more true positive interactions, which is crucial for DTI prediction.

\section{Comparison Against Literature Baseline (DeepCDA)}
\label{sec: eval_deepcda}
To contextualise the current results, they are compared against the DeepCDA model~\citep{Abbasi2020}, which is an established deep learning DTI prediction model that was evaluated on the BindingDB dataset.
Reported metrics are drawn from the comparative study of~\citep{Yang2024}.
\subsection{DeepCDA Architecture}
DeepCDA uses a dual-branch architecture where one branch processes protein sequences using an LSTM (Long Short-Term Memory) network, while the other branch processes drug SMILES strings using a CNN (Convolutional Neural Network).
The outputs of these two branches are then concatenated and passed through a fully connected (FC) classification layer to predict the binding interaction \citep{Abbasi2020}.
This architecture allows DeepCDA to capture both sequential dependencies in protein sequences and local patterns in drug SMILES, which are important for accurately predicting drug-target interactions.

This makes for a fair comparison with the model because both are binary classification models trained on BindingDB data, and both use deep learning architectures to process drug and target features separately before combining them for the final prediction.
Neither model is a direct implementation of the other, but they share enough similarities in their overall approach and training data to allow for a meaningful comparison of their performance metrics.

\subsection{Performance Comparison}
To establish a reference point before scaling up the data, the performance of the current model is evaluated against the reported metrics from DeepCDA, focusing on key metrics such as AUC-ROC, PR-AUC, precision, recall, and F1 score.

\subsubsection{Methodological Differences}
Before comparing the model's performance against DeepCDA, it's important to acknowledge the methodological differences between the two evaluations.
DeepCDA was evaluated using a 10-fold cross-validation approach, which involves partitioning the dataset into 10 subsets, training the model on 9 subsets and testing on the remaining subset, and repeating this process 10 times to obtain an average performance metric.
However, the model was evaluated using a fixed train/validation/test split, which means that the performance metrics we report are based on a single partitioning of the data.
This difference in evaluation methodology means that the metrics reported for DeepCDA are averaged across multiple runs and may be more robust to variations in the data, while these metrics are based on a single run and may be more sensitive to the specific train/test split used.
Therefore, while we can compare the reported metrics, we should do so with caution and acknowledge that the differences in evaluation methodology may impact the comparability of the results.

\subsubsection{Result}
Table~\ref{tab:deepcda_comparison} presents a comparison of this work against DeepCDA \citep{Abbasi2020}.
\begin{table}[h]
\centering
\caption{Performance comparison against literature baseline on BindingDB. DeepCDA results are reported from \citet{Yang2024}, evaluated using 10-fold cross-validation. This work uses a fixed train/validation/test split, so comparison is indicative rather than definitive.}
\label{tab:deepcda_comparison}
\resizebox{\textwidth}{!}{%
\begin{tabular}{lccccc}
\toprule
\textbf{Method} & \textbf{AUC-ROC} & \textbf{PR-AUC} & \textbf{F1}
& \textbf{Precision} & \textbf{Accuracy} \\
\midrule
DeepCDA \citep{Abbasi2020}    & 0.894 & 0.901 & 0.811 & 0.760 & 0.831 \\
\midrule
This work (Torch DNN, reduced dataset)         & 0.8670 & 0.8935 & 0.8302 & 0.8322 & 0.7963 \\
\bottomrule
\end{tabular}
}%
\end{table}

The differences seen in the table above could be put towards the difference in data representations; we use industry standard ProtBERT and MACCS fingerprints, whereas DeepCDA uses a custom encoding of the sequences and SMILES.
However, the biggest differences appear in how the models are evaluated, with DeepCDA using 10-fold cross-validation and this work using a fixed train/validation/test split.

It is worth noting that these results reflect an early checkpoint in the model development; subsequent sections detail further improvements through data.
A broader comparison against multiple literature baselines, including the improved final model, is presented in Section~\ref{sec:lit_comparison}.


\section{Data Sensitivity Analysis}
\label{sec: eval_data}
A data sensitivity analysis quantifies how training-set size and decision threshold affect model performance.
\subsection{Increased Sample Size}
After the move to the PyTorch framework, we have unlocked the ability to scale up the dataset to train the model on 50x more data.
PyTorch can handle the load better than a NumPy approach for the following reasons.

\begin{itemize}
    \item \textbf{GPU acceleration:} PyTorch offloads matrix multiplications to the GPU, NumPy runs everything on a single CPU.
    \item \textbf{Automatic batched computation:} PyTorch's DataLoader streams mini-batches to the GPU without loading the full dataset into working memory at once.
    The NumPy scripts manually slice arrays and do full matrix operations on CPU RAM.
    \item \textbf{Optimised autograd:} As mentioned earlier, autograd performs the backpropagation on a compiled C++ computation graph when NumPy computes gradients with python loops.
\end{itemize}

\begin{figure}[H]
    \centering
    \begin{minipage}{0.48\textwidth}
        \centering
        \includegraphics[width=\textwidth]{/Users/drs/Projects/DTI/Backend/app/Models/Torch/figures/dnn_adam_metrics_table_standalone.png}
        \caption*{Pytorch with Adam on reduced dataset}
    \end{minipage}
    \hfill
    \begin{minipage}{0.48\textwidth}
        \centering
        \includegraphics[width=\textwidth]{/Users/drs/Projects/DTI/Backend/app/Models/Experiments/grey_area/figures/original_metrics_table.png}
        \caption*{Pytorch with Adam on full dataset}
    \end{minipage}
    \caption{Analysis of increased sample size}
    \label{fig:full_dataset}
\end{figure}

The increased sample size has dramatically improved the performance of the model, being much more consistent over validation and test datasets.
The models were exactly the same, yet the change in data from ~18k to ~900k has reduced overfitting without redistributing the class balance, which we unravel in the following subsections.


    \subsection{Effect of Decision Threshold}
    \label{sec: data_threshold}
When evaluating the performance of a model, it's important to consider the thresholds used to determine the classification of predictions.
The model metrics are reported in Figure~\ref{fig:threshold_sweep} for different classification thresholds of 0.1, 0.3, 0.5, 0.7, and 0.9.
\begin{figure}[H]
    \centering
    \begin{minipage}{0.48\textwidth}
        \centering
        \includegraphics[width=\textwidth]{/Users/drs/Projects/DTI/Backend/app/Models/Torch/figures/dnn_adam_threshold_sweep.png}
        \caption*{Threshold Sweep Graph}
    \end{minipage}
    \hfill
    \begin{minipage}{0.48\textwidth}
        \centering
        \includegraphics[width=\textwidth]{/Users/drs/Projects/DTI/Backend/app/Models/Torch/figures/dnn_adam_threshold_sweep_table.png}
        \caption*{Threshold Sweep Metric Table}
    \end{minipage}
    \caption{Analysis of Threshold Levels on Model Performance.}
    \label{fig:threshold_sweep}
\end{figure}

After concluding the threshold sweep, we can identify an expected inverse relationship between Precision and Recall across all tested thresholds.
At a low threshold of 0.1, the model recorded a high recall of 0.945 but at the cost of lower precision at 0.750, suggesting that the model is classifying many instances as positive, which increases the likelihood of false positives.
As the threshold increases, precision improves to 0.897 and recall drops down to 0.621 at threshold 0.9, indicating that the model is now more conservative in its positive classifications, which reduces false positives but also increases false negatives.
The F1 score, which is the harmonic mean of precision and recall, peaks at a threshold of 0.3, suggesting that this is the optimal operating point for the model.
Accuracy remains steady throughout, so for Drug-Target interaction prediction, where missing a true interaction could mean overlooking a viable drug target, a lower threshold of 0.3 will be preferable.

\subsection{Grey Area Analysis}
\label{subsec:grey-area-analysis}
To further understand the impact of the \textbf{grey area} (weak predictions with no clear classification), adaptations to such thresholds will be tested to see how it affects model performance and the class imbalance issue.
Analysing the distribution of data, we can identify appropriate cutoffs to keep the distribution balanced between binders and non-binders using standard deviation.
The original thresholds of 5.30 and 7.00 were chosen to create a clear separation between binders and non-binders based off of literature standards, but this was not based off of the actual distribution of the data used in this work.

\begin{figure}[H]
    \centering
    \includegraphics[width=0.8\textwidth]{/Users/drs/Projects/DTI/Backend/app/Models/Experiments/grey_area/figures/paffinity_distribution.png}
    \caption{pAffinity distribution with median and mean marked}
    \label{fig:paffinity_distribution}
\end{figure}

\subsubsection{New Thresholds}
The original thresholds (5.30–7.00) were asymmetric with the upper bound being only 0.31 std above the mean, but the lower bound is 0.81 std below the mean.
Hence, the 61/39 imbalance within the original split dataset and the new approach of using more symmetric thresholds to give more naturally balanced classes with the following formula:
$\mu \pm k \cdot \sigma$
Where k is the number of standard deviations from the mean ($\mu$) ($\mu$ = 6.53, $\sigma$ = 1.52) and $\sigma$ is the standard deviation of the pAffinity distribution.
This ensures equal margins on both sides and produces the following new \textbf{grey area thresholds}:

\begin{itemize}
    \item \textbf{Original:} 5.30–7.00
    \item \textbf{k = 0.25:} 6.15–6.92
    \item \textbf{k = 0.5:} 5.77–7.30
    \item \textbf{No Grey Area:} median = 6.54
\end{itemize}

\subsubsection{Data Distribution Analysis}
After preprocessing the dataset again with the new thresholds, it created a more even class distribution, and can visualise the extent of the new grey areas in the Figure~\ref{fig:class_distribution}.
\begin{figure}[H]
    \centering
    \includegraphics[width=0.8\textwidth]{/Users/drs/Projects/DTI/Backend/app/Models/Experiments/grey_area/figures/class_distribution.png}
    \caption{Class distribution of new grey area thresholds compared to original thresholds}
    \label{fig:class_distribution}
\end{figure}

All the new grey area variations have much more even split between binders and non-binders, with no grey area giving the best balance 50.1/49.9 but will lose the ability to exclude weak predictions.

\begin{figure}[H]
    \centering
    \includegraphics[width=0.8\textwidth]{/Users/drs/Projects/DTI/Backend/app/Models/Experiments/grey_area/figures/paffinity_thresholds.png}
    \caption{pAffinity distribution with new grey area thresholds marked}
    \label{fig:new_thresholds_distribution}
\end{figure}

Figure~\ref{fig:new_thresholds_distribution} shows the new thresholds over the pAffinity distribution, with the original thresholds being more extreme and excluding a larger portion of the data as grey area, while the new thresholds are more central and include more data in the training set.

\subsubsection{Model Performance Analysis}
We ran the best model with each new threshold dataset to compare the performance of the model, here is what we discovered:

\begin{figure}[H]
    \centering
    \includegraphics[width=0.8\textwidth]{/Users/drs/Projects/DTI/Backend/app/Models/Experiments/grey_area/figures/comparison_loss.png}
    \caption{Comparison of training and validation loss between the thresholds ran on the same model}
    \label{fig:comparison_loss}
\end{figure}

\begin{figure}[H]
    \centering
    \includegraphics[width=0.8\textwidth]{/Users/drs/Projects/DTI/Backend/app/Models/Experiments/grey_area/figures/comparison_metrics.png}
    \caption{Comparison of metrics between the thresholds ran on the same model}
    \label{fig:comparison_metrics}
\end{figure}

Figures~\ref{fig:comparison_loss} and~\ref{fig:comparison_metrics} tell an interesting story, because despite the dataset being distributed more evenly with new grey area thresholds, the original threshold (5.3 < pAffinity < 7.0) with a ~60/40 split has performed the best.
Outperforming the other grey area thresholds across all metrics proves its validity and provides supporting evidence for removing the grey area as all the thresholds performed better than no threshold.


\section{Literature Comparison}
\label{sec:lit_comparison}
Having completed all iterations of model development, we now compare the final model against a broader set of deep learning DTI prediction methods evaluated on the BindingDB dataset.
Section~\ref{sec: eval_deepcda} presented a preliminary comparison against DeepCDA at an early stage of development; this comparison is extended to include TripletMulti-DTI, GAT, and MINDG, all reported in the comparative study of\citet{Yang2024}.

\subsection{Compared Architectures}
\label{subsec:compared-architectures}
Each of the compared methods takes a fundamentally different approach to modelling drug--target interactions, making this a diverse and informative benchmark.

\subsubsection{DeepCDA}
As discussed in Section~\ref{sec: eval_deepcda}, DeepCDA \citep{Abbasi2020} uses a dual-branch architecture with an LSTM for protein sequences and a CNN for drug SMILES strings, followed by a fully connected classification layer.

\subsubsection{TripletMulti-DTI}
TripletMulti-DTI employs a joint loss framework that jointly optimises a triplet loss alongside the binary classification objective.
It uses separate encoders for drugs and targets and learns a shared embedding space where interacting pairs are pulled closer together while non-interacting pairs are pushed apart \citep{DEHGHAN2023120754}.
This learning component gives the model a clearer representation of similarity between drug--target pairs beyond what a standard classification loss provides.

\subsubsection{GAT (Graph Attention Network)}
The GAT-based approach represents drugs as molecular graphs, where atoms are nodes and bonds are edges, and applies graph attention layers to learn molecular representations \citep{wang2021drugtargetinteractionpredictiongraph}.
Protein targets are encoded separately, and the combined representations are passed through layers for prediction.
Graph-based methods have the advantage of directly encoding molecular topology, but they suffer from being more computationally expensive.

\subsubsection{MINDG}
MINDG is the strongest baseline in the comparison, combining graph neural networks with multiscale interaction modelling \citep{Yang2024}.
It captures both local substructure patterns and global molecular properties and uses a dual-graph architecture to model drug and target features at multiple resolutions.
MINDG achieved the highest reported performance among the baselines across most metrics, making it the most challenging benchmark for this model.

\subsubsection{This Work}
This model uses a dual-branch fully connected DNN that processes pre-computed ProtBERT embeddings for protein targets and MACCS fingerprints for drug molecules.
Unlike the other methods, which learn feature representations end-to-end from raw sequences or molecular graphs, this approach relies on established, domain-specific featurisation pipelines.
The model was trained with a PyTorch model using the Adam optimiser with cosine annealing learning rate decay on the Ki/Kd subset of the full BindingDB dataset, as detailed in the preceding sections.

\subsection{Methodological Differences}
\label{subsec:methodological-differences}
As noted in Section~\ref{sec: eval_deepcda}, all four baseline methods were evaluated using 10-fold cross-validation, where the dataset is partitioned into 10 subsets and performance is averaged across 10 training runs.
The model was evaluated using a fixed train/validation/test split, meaning the reported metrics are based on a single partitioning of the data rather than an average over multiple.
Cross-validation generally produces more stable estimates with lower variance, and the reported standard deviations for the baselines confirm this.
The single-split evaluation may therefore be more sensitive to the specific data partition, and the absence of confidence intervals should be kept in mind when interpreting the results.
Despite this methodological difference, the comparison remains informative as all approaches are tested to unseen data drawn from the same source.

\subsection{Result}
Table~\ref{tab:performance-comparison} presents the full comparison.
All baseline results are reproduced from\citet{Yang2024}.

\begin{table}[htbp]
\centering
\caption{Performance comparison of different methods on BindingDB dataset.
Bold values indicate the best performance in each column.
Results for DeepCDA, TripletMulti-DTI, GAT, and MINDG are reproduced from \citet{Yang2024} and were evaluated using 10-fold cross-validation.
This work uses a fixed train/validation/test split.}
\label{tab:performance-comparison}
\resizebox{\textwidth}{!}{%
\begin{tabular}{lccccccc}
\toprule
\textbf{Method} & \textbf{PR-AUC} & \textbf{AUC-ROC} & \textbf{F1-score} & \textbf{Sensitivity} & \textbf{Specificity} & \textbf{Precision} & \textbf{Accuracy} \\
\midrule
DeepCDA          & $0.901 \pm 0.012$ & $0.894 \pm 0.003$ & $0.811 \pm 0.005$ & $0.872 \pm 0.013$ & $0.840 \pm 0.004$ & $0.760 \pm 0.016$ & $0.831 \pm 0.011$ \\
TripletMulti-DTI & $0.940 \pm 0.003$ & $0.931 \pm 0.001$ & $0.840 \pm 0.002$ & $0.917 \pm 0.022$ & $\mathbf{0.863 \pm 0.001}$ & $0.792 \pm 0.008$ & $0.865 \pm 0.005$ \\
GAT              & $0.923 \pm 0.002$ & $0.913 \pm 0.001$ & $0.755 \pm 0.021$ & $0.887 \pm 0.001$ & $0.851 \pm 0.001$ & $0.701 \pm 0.013$ & $0.775 \pm 0.012$ \\
MINDG            & $0.971 \pm 0.008$ & $0.951 \pm 0.004$ & $0.857 \pm 0.013$ & $0.923 \pm 0.006$ & $0.842 \pm 0.015$ & $0.800 \pm 0.013$ & $0.875 \pm 0.007$ \\
\midrule
Ours             & $\mathbf{0.9789}$ & $\mathbf{0.9703}$ & $\mathbf{0.9309}$ & $\mathbf{0.9283}$ & $0.9283$ & $\mathbf{0.9335}$ & $\mathbf{0.9162}$ \\
\bottomrule
\end{tabular}
}%
\end{table}

\subsection{Discussion}
Our model achieves the highest scores across six of the seven reported metrics, outperforming all four baselines in PR-AUC, AUC-ROC, F1-score, sensitivity, precision, and accuracy.
The margin of improvement over MINDG, the strongest baseline, is notable: the AUC-ROC of 0.9703 exceeds MINDG's 0.951, and the F1-score of 0.9309 represents a substantial improvement over MINDG's 0.857.
Precision sees the largest relative gain, with our model achieving 0.9335 compared to MINDG's 0.800, indicating that the model produces significantly fewer false positive predictions.

These results are particularly encouraging given the relative simplicity of our architecture.
While the baselines employ graph neural networks, attention mechanisms, contrastive learning, and multiscale interaction modelling, our model uses a straightforward fully connected DNN operating on pre-computed ProtBERT and MACCS features.
This suggests that the quality of input featurisation, specifically the use of ProtBERT embeddings and MACCS fingerprints, may be at least as important as architectural complexity for DTI prediction on this dataset.

However, the comparison should be interpreted with appropriate caution.
The use of a fixed split versus the baselines' 10-fold cross-validation means our reported metrics lack the variance estimates that would strengthen claims of statistical significance.
Not ignoring the caveats, the consistently strong performance across multiple metrics provides evidence that this approach is competitive with, and potentially exceeds, the current state of the art for DTI prediction on BindingDB, as reported in a 2024 study \citep{Yang2024}.
% ============================================================
\section{Evaluation Summary}
\label{sec: eval_summary}
After the full evaluation, we have determined the optimal model to be the Pytorch implementation with Adam optimiser on the full dataset, using a threshold of 0.3 for classification.
This model gave us the best overall performance with an F1 score of 0.9309 and an AUC-ROC of 0.9703, which is a significant improvement over the baseline model and is competitive with the literature baseline of DeepCDA and MINDG.
The most impactful improvement across iterations was the switch to the PyTorch framework, which provided a significant boost in performance because it unlocked the ability to use the full dataset instead of a subset.
The least impactful change for the model was the addition of a second hidden layer, which actually led to a drop in performance, likely due to overfitting given the limited data and lack of regularisation at that stage.
The main limitation with the evaluation is the differing methodologies used to evaluate our model and the literature baseline, with DeepCDA using 10-fold cross-validation and our model using a fixed train/validation/test split.
This means that the metrics we report for the model are based on a single run and may be more sensitive to the specific train/test split used, while the metrics reported for DeepCDA are averaged across multiple runs and may be more robust to variations in the data.